% !TeX program = xelatex
% 源：教学课件/A往年高一上学期课件/第4章/4.2.2指数函数的图象和性质课件.pptx
% 原稿含前置幂函数复习、4.2.1概念内容，均保留并分节。
\documentclass[aspectratio=169,10pt]{beamer}
\usepackage[UTF8,fontset=fandol]{ctex}
\IfFontExistsTF{Microsoft YaHei}{\setCJKsansfont{Microsoft YaHei}}{}
\usetheme{GaoZhongYuan}
% 教学内容层；原模板主题、图标及校园素材保持原样。
\usepackage{amsmath,amssymb,booktabs,array,pgfplots}
\pgfplotsset{compat=1.18}
\definecolor{lessonblue}{HTML}{225EA8}
\definecolor{lessonred}{HTML}{BA312D}
\definecolor{lessonpurple}{HTML}{82589F}
\setbeamerfont{normal text}{size=\large}
\setbeamerfont{block title}{size=\normalsize}
\setbeamerfont{block body}{size=\normalsize}
\setbeamertemplate{itemize item}{\color{campusgreen}\small$\bullet$}
\setbeamertemplate{enumerate item}{\color{campusgreen}\insertenumlabel.}
\setlength{\parskip}{.45em}
\setlength{\abovedisplayskip}{.5em}
\setlength{\belowdisplayskip}{.5em}
\newcommand{\R}{\mathbb R}
\newcommand{\Q}{\mathbb Q}
\newcommand{\ans}[1]{{\color{lessonred}#1}}
\newcommand{\key}[1]{{\color{campusgreen}\bfseries #1}}
\newcommand{\reason}[1]{{\color{lessonblue}#1}}
\newcommand{\gap}{\par\vspace{.6em}}
\pgfplotsset{lessonaxis/.style={axis lines=middle,axis line style={->},
 xlabel={$x$},ylabel={$y$},
 tick label style={font=\scriptsize},label style={font=\small},
 width=7.1cm,height=5.1cm,samples=150,clip=true,unbounded coords=jump,
 every axis plot/.append style={thick},
 legend style={font=\scriptsize,draw=none,fill=white,fill opacity=.9,text opacity=1},
 legend cell align=left}}
% 五个基本幂函数，均为可编辑矢量曲线。
\newcommand{\fivepowers}{%
\begin{tikzpicture}\begin{axis}[lessonaxis,xmin=-2.1,xmax=3.2,ymin=-2.1,ymax=3.6,
 xtick={-2,-1,1,2,3},ytick={-2,-1,1,2,3},legend pos=north west]
 \addplot[black,domain=-2:3.1]{x};\addlegendentry{$y=x$}
 \addplot[lessonblue,domain=-1.85:1.85]{x^2};\addlegendentry{$y=x^2$}
 \addplot[lessonred,domain=-1.3:1.5]{x^3};\addlegendentry{$y=x^3$}
 \addplot[campusgreen,domain=0:3.1]{sqrt(x)};\addlegendentry{$y=\sqrt{x}$}
 \addplot[lessonpurple,domain=.285:3.1]{1/x};\addlegendentry{$y=1/x$}
 \addplot[lessonpurple,domain=-2:-.48,forget plot]{1/x};
 \addplot[densely dashed,gray,forget plot] coordinates {(0,1)(1,1)(1,0)};
 \addplot[only marks,mark=*,mark size=1.3pt,forget plot] coordinates {(1,1)};
\end{axis}\end{tikzpicture}}
% 示意值只用于表达 0<q<m<1，未给原题补入数值条件。
\newcommand{\powerorder}{%
\begin{tikzpicture}\begin{axis}[lessonaxis,xmin=0,xmax=3.2,ymin=0,ymax=4,
 xtick={1,2,3},ytick={1,2,3},legend pos=north west]
 \addplot[lessonred,domain=0:1.95]{x^2};\addlegendentry{$y=x^p$}
 \addplot[campusgreen,domain=0:3.1]{x^.6};\addlegendentry{$y=x^m$}
 \addplot[lessonblue,domain=0:3.1]{x^.3};\addlegendentry{$y=x^q$}
 \addplot[lessonpurple,domain=.25:3.1]{1/x};\addlegendentry{$y=x^n$}
 \addplot[densely dashed,gray,forget plot] coordinates {(0,1)(1,1)(1,0)};
\end{axis}\end{tikzpicture}}
\newcommand{\expplot}[1][2]{%
\begin{tikzpicture}\begin{axis}[lessonaxis,xmin=-2.5,xmax=2.5,ymin=0,ymax=5,
 xtick={-2,-1,1,2},ytick={1,2,3,4},legend pos=north west]
 \addplot[lessonred,domain=-2.4:2.3]{2^x};\addlegendentry{$y=2^x$}
 \ifnum#1>1
 \addplot[lessonblue,domain=-2.3:2.4]{2^(-x)};\addlegendentry{$y=(1/2)^x$}
 \fi
 \ifnum#1>2
 \addplot[lessonpurple,domain=-2.4:1.45]{3^x};\addlegendentry{$y=3^x$}
 \addplot[campusgreen,domain=-1.45:2.4]{3^(-x)};\addlegendentry{$y=(1/3)^x$}
 \fi
 \addplot[only marks,mark=*,mark size=1.4pt,black,forget plot] coordinates {(0,1)};
\end{axis}\end{tikzpicture}}
% 四个选项依据原图特征重绘，B 为实立方根。
\newcommand{\rootchoices}{%
\foreach \letter/\kind in {A/1,B/2,C/3,D/4}{%
\begin{minipage}[t]{.24\textwidth}\centering
\begin{tikzpicture}\begin{axis}[lessonaxis,width=3.5cm,height=3.3cm,
 xmin=-1.6,xmax=1.7,ymin=-1.5,ymax=1.8,xtick={1},ytick={1},samples=201]
 \addplot[dashed,gray,domain=-1.4:1.6]{x};
 \ifnum\kind=1 \addplot[black,domain=-1.5:1.5]{tanh(x)};\fi
 \ifnum\kind=2 \addplot[black,domain=-1.5:1.5]{sign(x)*abs(x)^(1/3)};\fi
 \ifnum\kind=3 \addplot[black,domain=-1.12:1.2]{x^3};\fi
 \ifnum\kind=4 \addplot[black,domain=-1.5:1.5]{sign(x+.4)*abs(x+.4)^(1/3)};\fi
\end{axis}\end{tikzpicture}\\[-.5em]\letter
\end{minipage}}\par}

\title{4.2.2 指数函数的图象和性质}
\author{深圳中学高中园}
\date{}
\begin{document}
\campusmaketitle

% 原第1--5页，复习内容保留。
\begin{frame}[t]{课前复习：幂函数}
 \begin{itemize}
  \item 根据图象判断幂指数。
  \item 利用已知函数值确定幂函数。
  \item 利用幂函数的单调性比较大小。
 \end{itemize}
 \gap
 \key{幂函数中变量在底数位置，指数为常数。}
\end{frame}
% 原幂函数第18--20、23、25--27页；也用于指数函数原稿第2--5页复习。
\begin{frame}[t]{例2：立方根函数的图象}
 函数 $y=x^{1/3}=\sqrt[3]x$ 的图象是（\quad）。
 \par\rootchoices
 \pause
 \key{过 $(0,0)$、$(1,1)$，关于原点对称。}\\
 当 $x>1$ 时，$x^{1/3}<x$；当 $0<x<1$ 时，$x^{1/3}>x$。\\
 因而应选 \ans{B}。
\end{frame}

\begin{frame}[t]{例2：由图象判断指数大小}
 \begin{columns}[T,onlytextwidth]
  \begin{column}{.53\textwidth}\centering\powerorder\end{column}
  \begin{column}{.44\textwidth}
   根据图象中 $x>1$ 的部分，比较 $m,n,p,q$ 的大小。
   \pause
   \[n<0<q<m<1<p.\]
   因而
   \[\ans{n<q<m<p}.\]
   \reason{在 $x>1$ 时，指数越大，图象越高。}
  \end{column}
 \end{columns}
\end{frame}

\begin{frame}[t]{例3：待定指数求函数值}
 若 $f(x)$ 是幂函数，且 $f(8)=16$，求 $f(9),f(64)$。
 \pause
 设 $f(x)=x^\alpha$。由 $8^\alpha=16$ 得
 \[2^{3\alpha}=2^4\quad\Longrightarrow\quad\alpha=\frac43.\]
 因此 $f(x)=x^{4/3}$，
 \[\ans{f(9)=9^{4/3}=9\sqrt[3]9},\qquad
   \ans{f(64)=(2^6)^{4/3}=256}.\]
 \reason{函数类型已知时，先设标准形式，再求参数。}
\end{frame}

\begin{frame}[t]{例4：同指数幂的大小}
 比较各组中两个数的大小：
 \[(1)\ 1.5^{1.4},\ 1.6^{1.4};\qquad
   (2)\ 1.5^{0.4},\ 1.6^{0.4}.\]
 \pause
 \begin{itemize}
  \item $y=x^{1.4}$ 在 $[0,+\infty)$ 上递增，故
   \[\ans{1.5^{1.4}<1.6^{1.4}}.\]
  \item $y=x^{0.4}$ 在 $[0,+\infty)$ 上递增，故
   \[\ans{1.5^{0.4}<1.6^{0.4}}.\]
 \end{itemize}
\end{frame}

\begin{frame}[t]{例4：负指数与比较方法}
 \[(3)\ 1.5^{-1.5},\ 1.6^{-1.5}.\]
 \pause
 $y=x^{-1.5}$ 在 $(0,+\infty)$ 上递减，且 $1.5<1.6$，故
 \[\ans{1.5^{-1.5}>1.6^{-1.5}}.\]
 \gap
 \key{用单调性直接比较时，两个自变量应处于同一单调区间。}
 \gap
 跨区间时应另作分析，例如：$-1<1$，但 $(-1)^{-1}<1^{-1}$。\\
 $y=1/x$ 虽在两侧分别递减，却不在整个定义域上递减。
\end{frame}


% 原第6--8页
\begin{frame}[t]{任务一：指数函数的概念}
 在幂的表达式 $a^x$ 中，若底数 $a$ 为常数，指数 $x$ 为自变量，\\
 就得到函数
 \[\boxed{y=a^x\qquad(a>0\text{ 且 }a\ne1)}.\]
 它叫做\key{指数函数}，定义域为 $\R$。
 \pause
 \gap
 \begin{columns}[T,onlytextwidth]
  \begin{column}{.46\textwidth}
   \key{幂函数 $y=x^\alpha$}\\变量在底数，指数固定。
  \end{column}
  \begin{column}{.48\textwidth}
   \key{指数函数 $y=a^x$}\\底数固定，变量在指数。
  \end{column}
 \end{columns}
\end{frame}

\begin{frame}[t]{任务一：底数为何受限制}
 \begin{itemize}
  \item 若 $a=0$，$x\le0$ 时 $0^x$ 无定义，不能在 $\R$ 上定义。
  \item 若 $a<0$，如 $(-4)^{1/2}$ 在实数范围内无意义，\\
        因而也不能对每个实数 $x$ 定义 $a^x$。
  \item 若 $a=1$，则 $1^x=1$，得到常函数。
 \end{itemize}
 \pause
 \gap
 因此，指数函数的定义规定 $\ans{a>0\text{ 且 }a\ne1}$。
\end{frame}

% 原第9页，判断题与求值分开。
\begin{frame}[t]{概念自检：判断正误}
 判断下列说法是否正确：
 \begin{enumerate}
  \item $y=x^2$ 是指数函数。
  \item 指数函数 $y=a^x$ 的底数 $a$ 可以为负数。
  \item $y=2^x-1$ 是指数函数。
 \end{enumerate}
 \pause
 \ans{三项均错误。}
 \gap
 第1项为幂函数；第2项违反底数条件；\\
 第3项是由指数函数平移得到的函数，不符合定义的标准形式。
\end{frame}

\begin{frame}[t]{概念自检：由函数值求底数}
 若 $f(x)$ 是指数函数，且 $f(2)=2$，求 $f(x)$。
 \pause
 设 $f(x)=a^x\ (a>0,a\ne1)$，由
 \[a^2=2\]
 及 $a>0$，得 $a=\sqrt2$。所以
 \[\ans{f(x)=(\sqrt2)^x}.\]
\end{frame}

% 原第10页
\begin{frame}[t]{例1：识别指数函数}
 下列函数中，哪些是指数函数？
 \[\text{(1) }y=2(\sqrt2)^x;\qquad
   \text{(2) }y=2^{x-1};\qquad
   \text{(3) }y=\left(\frac\pi2\right)^x.\]
 \pause
 \begin{itemize}
  \item (1) 的系数为 $2$，不符合 $y=a^x$。
  \item (2) 可写成 $y=\frac12\cdot2^x$，也不符合标准形式。
  \item (3) 的底数 $\pi/2>0$ 且不等于 $1$，符合定义。
 \end{itemize}
 \ans{答案：(3)。}
\end{frame}

% 原第11页
\begin{frame}[t]{例1：由定义求参数}
 若 $y=(k+2)a^x+2-b\ (a>0,a\ne1)$ 是指数函数，求 $k,b$。
 \pause
 比较标准形式 $y=a^x$，得
 \[\begin{cases}k+2=1,\\2-b=0.\end{cases}\]
 所以
 \[\ans{k=-1,\qquad b=2}.\]
\end{frame}

% 原第12页
\begin{frame}[t]{例1：由已知点确定函数}
 若 $f(x)$ 是指数函数，且 $f(2)=9$，求 $f(x)$。
 \pause
 设 $f(x)=a^x\ (a>0,a\ne1)$，则
 \[a^2=9.\]
 由 $a>0$ 得 $a=3$，所以
 \[\ans{f(x)=3^x}.\]
\end{frame}

% 原第13页
\begin{frame}[t]{跟踪训练：底数条件}
 若 $y=(2a-1)^x$ 是指数函数，则 $a$ 的取值范围是（\quad）。
 \normalsize
 \begin{align*}
  \mathrm A.&\ (0,1)\cup(1,+\infty);&
  \mathrm B.&\ [0,1)\cup(1,+\infty);\\[.5em]
  \mathrm C.&\ (\tfrac12,1)\cup(1,+\infty);&
  \mathrm D.&\ [\tfrac12,+\infty).
 \end{align*}
 \pause
 \[2a-1>0,\quad 2a-1\ne1\quad\Longrightarrow\quad
   a>\frac12,\quad a\ne1.\]
 \ans{答案：C。}
\end{frame}

% 原第14--15页
\begin{frame}[t]{任务二：指数增长与衰减}
 \begin{columns}[T,onlytextwidth]
  \begin{column}{.47\textwidth}
   \key{当 $a>1$ 时}\\
   $a^x$ 随 $x$ 增大而增大。\\
   $x$ 充分增大时，函数值可任意大。\\
   这是指数增长，也称指数爆炸。
  \end{column}
  \begin{column}{.47\textwidth}
   \key{当 $0<a<1$ 时}\\
   $a^x$ 随 $x$ 增大而减小。\\
   $x$ 充分增大时，函数值无限接近 $0$。\\
   这是指数衰减。
  \end{column}
 \end{columns}
 \pause
 \gap
 \[2^0=1,\quad2^5=32,\quad2^{10}=1024;\qquad
  (\tfrac12)^0=1,\quad(\tfrac12)^5=\tfrac1{32}.\]
\end{frame}

% 原第16--17页
\begin{frame}[t]{任务三：描点作出指数图象}
 用描点法作 $y=2^x$ 的图象：列表、描点、用平滑曲线连接。
 \begin{columns}[T,onlytextwidth]
  \begin{column}{.38\textwidth}
   \centering\normalsize\renewcommand{\arraystretch}{1.55}
   \begin{tabular}{c|c}
    $x$ & $2^x$\\\hline
    $-2$ & $\frac14$\\$-1$ & $\frac12$\\$0$ & $1$\\$1$ & $2$\\$2$ & $4$
   \end{tabular}
   \pause
   \par\gap
   曲线经过 $(0,1)$，\\始终位于 $x$ 轴上方。
  \end{column}
  \begin{column}{.58\textwidth}\centering\expplot[1]\end{column}
 \end{columns}
\end{frame}

% 原第18页
\begin{frame}[t]{任务三：衰减函数的图象}
 再作 $y=(1/2)^x$ 的图象，与 $y=2^x$ 比较。
 \begin{columns}[T,onlytextwidth]
  \begin{column}{.38\textwidth}
   \centering\normalsize\renewcommand{\arraystretch}{1.55}
   \begin{tabular}{c|c|c}
    $x$ & $2^x$ & $(1/2)^x$\\\hline
    $-2$ & $\frac14$ & $4$\\$-1$ & $\frac12$ & $2$\\
    $0$ & $1$ & $1$\\$1$ & $2$ & $\frac12$\\$2$ & $4$ & $\frac14$
   \end{tabular}
  \end{column}
  \begin{column}{.58\textwidth}\centering\expplot\end{column}
 \end{columns}
\end{frame}

% 原第19--20页
\begin{frame}[t]{任务四：互为倒数的底数}
 \begin{columns}[T,onlytextwidth]
  \begin{column}{.52\textwidth}\centering\expplot\end{column}
  \begin{column}{.45\textwidth}
   设 $f(x)=a^x$，则
   \[\left(\frac1a\right)^x=a^{-x}=f(-x).\]
   \pause
   因而 $y=a^x$ 与 $y=(1/a)^x$ 的图象\key{关于 $y$ 轴对称}。
   \gap
   两者都经过 $(0,1)$。
  \end{column}
 \end{columns}
\end{frame}

% 原第20--22页
\begin{frame}[t]{任务四：不同底数的图象}
 \begin{columns}[T,onlytextwidth]
  \begin{column}{.52\textwidth}\centering\expplot[4]\end{column}
  \begin{column}{.45\textwidth}
   观察 $2^x,3^x,(1/2)^x,(1/3)^x$：
   \begin{enumerate}
    \item 图象经过哪些象限？
    \item 共同经过哪个点？
    \item 增减性与底数有何关系？
   \end{enumerate}
   \pause
   $a>1$ 时递增；\\$0<a<1$ 时递减。
  \end{column}
 \end{columns}
\end{frame}

% 原第23页，补充x<0一行。
\begin{frame}[t]{任务五：指数函数的性质}
 \centering\normalsize\renewcommand{\arraystretch}{1.65}
 \begin{tabular}{c|c|c}
  \toprule $y=a^x$ & $a>1$ & $0<a<1$\\\midrule
  定义域 & $\R$ & $\R$\\
  值域 & $(0,+\infty)$ & $(0,+\infty)$\\
  定点 & $(0,1)$ & $(0,1)$\\
  $x>0$ & $a^x>1$ & $0<a^x<1$\\
  $x<0$ & $0<a^x<1$ & $a^x>1$\\
  单调性 & 在 $\R$ 上递增 & 在 $\R$ 上递减\\\bottomrule
 \end{tabular}
 \gap
 \key{图象位于第一、二象限，并经过 $y$ 轴上的点 $(0,1)$。}
\end{frame}

% 原第24页
\begin{frame}[t]{例2：指数式大小比较}
 比较各组中两个数的大小：
 \begin{align*}
  (1)&\quad 1.7^{2.5}\quad\text{与}\quad1.7^3;\\[.5em]
  (2)&\quad 0.8^{-\sqrt2}\quad\text{与}\quad0.8^{-\sqrt3};\\[.5em]
  (3)&\quad 1.7^{0.3}\quad\text{与}\quad0.9^{3.1}.
 \end{align*}
 \pause
 \normalsize
 (1) $1.7>1$，递增：$\ans{1.7^{2.5}<1.7^3}$。\\[.3em]
 (2) $0<0.8<1$，递减；$-\sqrt2>-\sqrt3$：$\ans{0.8^{-\sqrt2}<0.8^{-\sqrt3}}$。\\[.3em]
 (3) 借助 $1$：$\ans{1.7^{0.3}>1>0.9^{3.1}}$。
\end{frame}

% 原第25页
\begin{frame}[t]{跟踪训练：指数式比较}
 比较各组中两个数的大小：
 \begin{align*}
  (1)&\quad 0.99^{3.3}\quad\text{与}\quad0.99^{4.5};\\[.5em]
  (2)&\quad 1.01^{2.7}\quad\text{与}\quad1.01^{3.5};\\[.5em]
  (3)&\quad 2.3^{-0.5}\quad\text{与}\quad0.2^{-0.1}.
 \end{align*}
 \pause
 \[\ans{0.99^{3.3}>0.99^{4.5}},\qquad
   \ans{1.01^{2.7}<1.01^{3.5}}.\]
 \[\ans{2.3^{-0.5}<1<0.2^{-0.1}}.\]
\end{frame}

% 原第26页；补充同指数类型以串联两课。
\begin{frame}[t]{大小比较的方法}
 \begin{itemize}
  \item \key{底数相同，指数不同}：利用指数函数的单调性。
  \item \key{指数相同，底数不同}：在合适区间利用幂函数的单调性。
  \item \key{底数、指数均不同}：尝试化为同底数、同指数，\\
        或借助 $1$ 等中间量比较。
 \end{itemize}
 \gap
 \reason{比较前，先看清哪个量相同、哪个量在变化。}
\end{frame}

% 原第27页
\begin{frame}[t]{例3：指数型函数的定点}
 函数 $y=a^{x-1}\ (a>0,a\ne1)$ 的图象经过什么定点？
 \pause
 令 $x-1=0$，得 $x=1$，此时 $y=a^0=1$。
 \[\ans{\text{定点为 }(1,1)}.\]
 \pause
 \key{跟踪训练}\quad $y=2a^{x+1}-3\ (a>0,a\ne1)$ 呢？
 \pause
 令 $x+1=0$，得 $x=-1$，$y=2\cdot1-3=-1$。
 \[\ans{\text{定点为 }(-1,-1)}.\]
\end{frame}

% 原第28页，补全D项的底数约束。
\begin{frame}[t]{当堂检测：性质辨析}
 （多选）下列判断正确的是（\quad）。
 \normalsize
 \begin{enumerate}[A.]
  \item $a^3>a^2\ (a>0\text{ 且 }a\ne1)$。
  \item $y=(1/2)^x\ (x\ge0)$ 的值域为 $(0,1]$。
  \item $y=2^{1/x}$ 的定义域为 $\R$。
  \item $y=a^{x-1}-1\ (a>0\text{ 且 }a\ne1)$ 的图象过定点 $(1,0)$。
 \end{enumerate}
 \pause
 \ans{答案：B、D。}\\
 A 在 $0<a<1$ 时不成立；C 应排除 $x=0$。\\
 B 的函数值可趋近 $0$ 但取不到 $0$；D 中令 $x=1$ 即得 $y=0$。
\end{frame}

% 原第29页
\begin{frame}[t]{课堂小结}
 \begin{itemize}
  \item 指数函数 $y=a^x$：$a>0$ 且 $a\ne1$，定义域为 $\R$。
  \item 值域为 $(0,+\infty)$，图象经过定点 $(0,1)$。
  \item $a>1$ 时递增，$0<a<1$ 时递减。
  \item 应用：比较大小，求指数型函数的定点。
 \end{itemize}
 \gap
 \key{数形结合、归纳类比、分类讨论。}
\end{frame}

% 原第30页；将有限小数前的等号纠正为约等号。
\begin{frame}[t]{课后作业与思考}
 \key{课后作业}
 \begin{itemize}
  \item 书本第113页练习第3题、习题4.2第5题。
  \item 同步训练第30页。
 \end{itemize}
 \gap
 \[1.01^{365}\approx37.8,\qquad0.99^{365}\approx0.03.\]
 \normalsize
 勤学如初见之苗，不见其增，日有所长；\\
 辍学如磨刀之石，不见其损，日有所亏。
\end{frame}
\campusclosing[谢谢]
\end{document}
